\part{Vorhersage und Politik}

\chapter{Was ein gutes Modell nicht kann}

\section{Zwei Zwecke}

\begin{merksatz}
\boxtitel{Vorhersagen und Eingreifen}
Ein Modell der Arztbesuche kann zweierlei leisten:

\begin{center}
{\small
\renewcommand{\arraystretch}{1.55}
\begin{tabular}{@{}P{2.4cm}P{6.2cm}@{}}
\toprule
\textbf{Vorhersage} & Wie viele Besuche nächstes Jahr --- wenn alles so weitergeht?\\
\textbf{Politik} & Wie viele Besuche --- wenn man die Zuzahlung ändert?\\
\bottomrule
\end{tabular}}
\end{center}

\begin{center}
\begin{minipage}{0.9\textwidth}
\textbf{Band 22 behandelte die erste Frage, Band 23 die zweite. Dieses Kapitel zeigt, dass ein Modell die erste glänzend beantworten und bei der zweiten vollständig versagen kann --- und zwar aus einem Grund, der nicht in der Statistik liegt.}
\end{minipage}
\end{center}
\end{merksatz}

\section{Die Kritik von Lucas}

\begin{merksatz}
\boxtitel{Ein Modell mit Erwartungen}
Versicherte reagieren nicht nur auf die heutige Zuzahlung, sondern auf die erwartete:
\[ y_t = 8 + \beta\,z_t + \gamma\,(z^{e}_{t+1} - z_t) + \varepsilon_t , \qquad \beta = -0{,}10,\quad \gamma = 0{,}05 \]
--- wer eine Erhöhung erwartet, zieht Besuche vor.

\schritt{Wie die Erwartung entsteht}
Die Zuzahlung folgt einer Regel der Kasse:
\[ z_{t+1} = 10 + \phi\,(z_t - 10) + \eta_{t+1}, \qquad z^{e}_{t+1} = 10 + \phi\,(z_t-10) \]

\schritt{Was man schätzt}
Regressiert man $y$ auf $z$, so erhält man nicht $\beta$, sondern
\[ \beta + \gamma\,(\phi - 1) . \]

\begin{center}
{\small
\renewcommand{\arraystretch}{1.48}
\begin{tabular}{@{}P{2.8cm}P{1.4cm}P{2.0cm}P{2.0cm}@{}}
\toprule
\textbf{Regel} & $\phi$ & \textbf{Formel} & \textbf{simuliert}\\
\midrule
träge Regel & $0{,}8$ & $-0{,}110$ & $-0{,}103$\\
sofortige Rückkehr & $0{,}0$ & $-0{,}150$ & $-0{,}141$\\
\bottomrule
\end{tabular}}
\end{center}

\begin{center}
\begin{minipage}{0.9\textwidth}
\textbf{Der geschätzte Koeffizient hängt von der Regel ab, unter der die Daten entstanden. Ändert die Kasse die Regel, ändert sich der Koeffizient --- obwohl sich am Verhalten der Versicherten nichts geändert hat.}
\end{minipage}
\end{center}
\end{merksatz}

\begin{merksatz}
\boxtitel{Satz 13.1 --- die Kritik von Lucas}
Die Koeffizienten eines Modells, das Erwartungen nicht ausdrücklich enthält, sind Funktionen der Politikregel. Sie sind daher für die Bewertung einer \kursiv{anderen} Politikregel ungeeignet.

\schritt{Der Gedanke}
Die \Begriff{tiefen Parameter} --- hier $\beta$ und $\gamma$ --- beschreiben Vorlieben und Technik und ändern sich mit der Politik nicht. Die \Begriff{reduzierte Form} --- hier $\beta+\gamma(\phi-1)$ --- verknüpft sie mit der Regel $\phi$. Wer die reduzierte Form schätzt und damit eine neue Regel $\phi'$ bewertet, benutzt den falschen Koeffizienten; der Fehler ist $\gamma(\phi'-\phi)$. $\blacksquare$

\schritt{Im Beispiel}
Eine Kasse wechselt von der trägen Regel zur sofortigen Rückkehr. Das alte Modell sagt für eine Erhöhung um zehn Euro $-1{,}03$ Besuche vorher; tatsächlich sind es $-1{,}41$ --- ein Fehler von fast vierzig Prozent, und das bei einem Modell, das unter der alten Regel vollkommen stimmte.

\begin{center}
\begin{minipage}{0.88\textwidth}
\kursiv{\textbf{Die Kritik trifft nicht die Statistik, sondern den Gegenstand}: Menschen reagieren auf Regeln, und ein Modell, das die Reaktion auf die alte Regel beschreibt, weiß nichts über die Reaktion auf eine neue. Band 23 kannte das Problem als Übertragbarkeit --- dort ging es um andere Populationen, hier um dieselbe Population unter anderen Bedingungen. Das zweite ist schwerer, weil man die Bedingungen nicht beobachten kann, bevor man sie herstellt.}
\end{minipage}
\end{center}
\end{merksatz}

\section{Der Vorhersagewettbewerb}

\begin{merksatz}
\boxtitel{Wer am besten vorhersagt}
Eine Reihe, erzeugt als AR(2), zweihundert Punkte zum Schätzen, hundert zum Prüfen. Vorhersage jeweils einen Schritt voraus, mit fortlaufender Neuschätzung:

\begin{center}
{\small
\renewcommand{\arraystretch}{1.48}
\begin{tabular}{@{}P{3.2cm}P{2.4cm}P{2.4cm}@{}}
\toprule
\textbf{Modell} & \textbf{Fehler außerhalb} & \textbf{Fehler innerhalb}\\
\midrule
Mittelwert & $1{,}760$ & \\
letzter Wert & $1{,}056$ & \\
AR(1) & $\mathbf{1{,}011}$ & \\
AR(2), wahre Form & $1{,}015$ & $1{,}000$\\
AR(12) & $1{,}050$ & \\
AR(24) & $1{,}066$ & $\mathbf{0{,}960}$\\
Mittel aus 1, 2, 12 & $1{,}022$ & \\
\midrule
wahre Koeffizienten & $1{,}025$ & \\
\bottomrule
\end{tabular}}
\end{center}

\begin{center}
\begin{minipage}{0.9\textwidth}
\textbf{Das Modell mit der besten Anpassung ist das schlechteste in der Vorhersage. Das falsche, zu kleine Modell schlägt sogar die wahre Form.}
\end{minipage}
\end{center}

\begin{center}
\begin{minipage}{0.88\textwidth}
\kursiv{\textbf{Das ist die Bias-Varianz-Abwägung aus Band 22, Kapitel~2}: Das AR(1) hat eine Verzerrung --- es lässt den zweiten Koeffizienten weg --- und gewinnt dafür an Stabilität der Schätzung. Bei zweihundert Beobachtungen wiegt das schwerer. Der Fehler des wahren Modells mit geschätzten Koeffizienten liegt über dem mit wahren Koeffizienten; die Differenz ist der Preis des Schätzens, und ein kleineres Modell zahlt weniger davon.}
\end{minipage}
\end{center}
\end{merksatz}

\begin{merksatz}
\boxtitel{Vier Lehren aus fünfzig Jahren Wettbewerben}
\begin{center}
{\small
\renewcommand{\arraystretch}{1.52}
\begin{tabular}{@{}P{1.0cm}P{7.4cm}@{}}
\toprule
$1$ & Einfache Modelle schlagen komplizierte --- regelmäßig und über Jahrzehnte.\\
$2$ & Kombinationen schlagen jedes Einzelmodell --- fast immer.\\
$3$ & Die Anpassung an die Vergangenheit sagt wenig über die Vorhersagegüte.\\
$4$ & Der Vorsprung guter Modelle vor dem naiven ist klein.\\
\bottomrule
\end{tabular}}
\end{center}

\begin{center}
\begin{minipage}{0.88\textwidth}
\kursiv{\textbf{Die vierte Zeile ist ernüchternd und wichtig}: Zwischen \enquote{letzter Wert} und bestem Modell liegen in der Tabelle vier Prozent. Bei Wechselkursen ist es seit 1983 bekannt, dass gar kein Modell die Irrfahrt schlägt --- das Rätsel von Meese und Rogoff. Wer große Vorhersagegewinne verspricht, hat sie meist innerhalb der Stichprobe gemessen.}
\end{minipage}
\end{center}
\end{merksatz}

\section{Der Bruchfall}

\begin{warnung}
\boxtitelr{Wenn das beste Vorhersagemodell die falsche Politik empfiehlt}
\begin{center}
{\small
\renewcommand{\arraystretch}{1.52}
\begin{tabular}{@{}P{2.8cm}P{5.8cm}@{}}
\toprule
\textbf{Vorhersage} & braucht Korrelationen, die stabil bleiben, solange sich nichts ändert.\\
\textbf{Politik} & braucht Mechanismen, die stabil bleiben, wenn man eingreift.\\
\textbf{Das Dilemma} & Die Modelle, die das eine am besten können, sind oft die, die das andere gar nicht können --- weil sie die Struktur, die bei Eingriffen zählt, zugunsten der Anpassung weglassen.\\
\bottomrule
\end{tabular}}
\end{center}

\begin{center}
\begin{minipage}{0.88\textwidth}
\kursiv{\textbf{Ein Beispiel aus der Klinik}: Das beste Vorhersagemodell für Wiederaufnahmen enthält die Zahl der bisherigen Aufenthalte als stärksten Regressor. Es sagt hervorragend vorher und ist für die Frage, was eine Wiederaufnahme \kursiv{verhindert}, wertlos --- die Aufenthalte sind Symptom, nicht Hebel. \textbf{Der Streit zwischen Vorhersage und Struktur ist der Streit zwischen Band 22 und Band 23, und die Ökonometrie hat ihn seit 1976 in aller Schärfe ausgetragen.}}
\end{minipage}
\end{center}
\end{warnung}

\section{Übungsaufgaben}

\begin{aufgabe}
\boxtitel{Aufgabe 13.1 --- Die reduzierte Form herleiten}
Zeigen Sie, dass die Regression von $y$ auf $z$ im Erwartungsmodell den Koeffizienten $\beta + \gamma(\phi-1)$ liefert.

\loesung
\schritt{Schritt 1 --- Die Erwartung einsetzen}
\[ z^{e}_{t+1} - z_t = 10 + \phi(z_t-10) - z_t = (\phi-1)(z_t-10) \]

\schritt{Schritt 2 --- In die Verhaltensgleichung}
\[ y_t = 8 + \beta z_t + \gamma(\phi-1)(z_t-10) + \varepsilon_t = \bigl[8-10\gamma(\phi-1)\bigr] + \bigl[\beta+\gamma(\phi-1)\bigr]z_t + \varepsilon_t \]

\schritt{Schritt 3 --- Die Probe}
$\phi = 0{,}8$: $-0{,}10 + 0{,}05\cdot(-0{,}2) = -0{,}110$; simuliert $-0{,}103$. \checkmark

\schritt{Schritt 4 --- Was identifiziert ist}
\begin{center}
\begin{minipage}{0.9\textwidth}
\textbf{Aus Daten unter einer einzigen Regel sind $\beta$ und $\gamma$ nicht trennbar --- nur ihre Kombination. Erst Daten unter zwei Regeln identifizieren beide.}
\end{minipage}
\end{center}

\begin{center}
\begin{minipage}{0.88\textwidth}
\kursiv{\textbf{Das ist Kapitel~2 in zeitlicher Gestalt}: Ein Parameter, zwei Unbekannte. Und es zeigt, worauf die Antwort auf Lucas hinausläuft: Man braucht Regelwechsel in den Daten --- natürliche Versuche, Reformen, Regimeänderungen ---, um die tiefen Parameter zu bestimmen. Die Glaubwürdigkeitsrevolution aus Kapitel~1 und die Lucas-Kritik sind insofern Verbündete, auch wenn ihre Vertreter sich selten so sehen.}
\end{minipage}
\end{center}
\end{aufgabe}

\begin{aufgabe}
\boxtitel{Aufgabe 13.2 --- Vorhersagen vergleichen}
Zwei Modelle haben über hundert Prüfpunkte quadratische Fehler mit Mittelwert $1{,}02$ und $1{,}08$. Ist das zweite schlechter?

\loesung
\schritt{Schritt 1 --- Die Frage richtig stellen}
Die Differenz der Fehlerquadrate je Zeitpunkt, $d_t = e_{1t}^{2} - e_{2t}^{2}$, hat Mittelwert $-0{,}06$. Ist das von null verschieden?

\schritt{Schritt 2 --- Der Standardfehler}
Die $d_t$ sind autokorreliert --- Satz 10.4 gilt. Mit einer Streuung der $d_t$ von etwa $1{,}5$ und wirksamem Umfang $50$:
\[ \text{Standardfehler} \approx \frac{1{,}5}{\sqrt{50}} = 0{,}21, \qquad t = \frac{-0{,}06}{0{,}21} = -0{,}29 \]

\begin{center}
\begin{minipage}{0.9\textwidth}
\textbf{Der Unterschied ist weit von jeder Signifikanz entfernt. Hundert Prüfpunkte können Modelle, die sich um sechs Prozent unterscheiden, nicht trennen.}
\end{minipage}
\end{center}

\schritt{Schritt 3 --- Die Einordnung}
\begin{center}
\begin{minipage}{0.88\textwidth}
\kursiv{\textbf{Das ist der Test von Diebold und Mariano}, und sein häufigstes Ergebnis ist: nicht unterscheidbar. Das erklärt die zweite Lehre aus den Wettbewerben --- wenn man nicht weiß, welches Modell besser ist, ist das Mittel aus beiden eine Absicherung. Und es erklärt, warum die Tabelle in Abschnitt 13.3 mit Vorsicht zu lesen ist: Die Rangfolge der AR-Modelle liegt innerhalb der Unsicherheit.}
\end{minipage}
\end{center}
\end{aufgabe}

\begin{aufgabe}
\boxtitel{Aufgabe 13.3 --- Wann die Kritik beißt}
Für welche der folgenden Fragen ist die Lucas-Kritik ernst, für welche nicht?
\begin{center}
{\small
\renewcommand{\arraystretch}{1.5}
\begin{tabular}{@{}P{1.0cm}P{7.4cm}@{}}
\toprule
(a) & Wirkung einer einmaligen, unangekündigten Zuzahlungserhöhung.\\
(b) & Wirkung einer dauerhaften Regeländerung, die Zuzahlungen an Vorjahresausgaben koppelt.\\
(c) & Wirkung einer Grippeimpfung auf Arztbesuche.\\
\bottomrule
\end{tabular}}
\end{center}

\loesung
\schritt{(a)}
\begin{center}
\begin{minipage}{0.9\textwidth}
\textbf{Kaum.} Eine einmalige Überraschung ändert keine Erwartungen; das geschätzte $\beta$ gilt.
\end{minipage}
\end{center}

\schritt{(b)}
\begin{center}
\begin{minipage}{0.9\textwidth}
\textbf{Voll.} Die Versicherten lernen die Regel, bilden Erwartungen über künftige Zuzahlungen und ändern ihr Verhalten --- genau der Fall aus Abschnitt 13.2.
\end{minipage}
\end{center}

\schritt{(c)}
\begin{center}
\begin{minipage}{0.9\textwidth}
\textbf{Nicht.} Die Wirkung einer Impfung ist Biologie, keine Reaktion auf Regeln. Die Kritik betrifft Verhaltensparameter, nicht technische.
\end{minipage}
\end{center}

\schritt{Die Regel}
\begin{center}
\begin{minipage}{0.88\textwidth}
\kursiv{\textbf{Die Kritik beißt, wo Menschen auf die Politik reagieren, deren Wirkung man schätzen will} --- und sie beißt umso stärker, je besser die Politik vorhersehbar ist. Das hat eine praktische Folge: Wer eine Maßnahme evaluiert, sollte fragen, ob die Betroffenen sie kommen sahen. Eine Reform, die ein Jahr vorher angekündigt wurde, hat ihre Wirkung teilweise schon vor dem Stichtag entfaltet --- und die doppelte Differenz aus Kapitel~9 misst dann zu wenig.}
\end{minipage}
\end{center}
\end{aufgabe}

\begin{geschichte}
\boxtitelg{Ein Aufsatz und seine Folgen}
\Person{Robert Lucas} trug die Kritik 1973 auf einer Konferenz vor und veröffentlichte sie 1976 unter dem Titel \kursiv{Econometric Policy Evaluation: A Critique}.

\begin{center}
\begin{minipage}{0.9\textwidth}
\kursiv{Der Anlass war das Versagen der großen makroökonomischen Modelle in den 1970er Jahren: Sie hatten eine stabile Austauschbeziehung zwischen Inflation und Arbeitslosigkeit geschätzt, und als die Politik sie ausnutzte, verschwand sie. \textbf{Lucas' Erklärung: Die Beziehung war eine reduzierte Form, und die Erwartungen hatten sich der Politik angepasst.} Der Nobelpreis kam 1995.}
\end{minipage}
\end{center}

\schritt{Die Antwort}
\begin{center}
\begin{minipage}{0.9\textwidth}
\kursiv{Die Makroökonomie reagierte mit Modellen, die Erwartungen ausdrücklich enthalten und nur tiefe Parameter schätzen --- um den Preis, dass diese Modelle viel annehmen und wenig vorhersagen. \textbf{\Person{Sims} antwortete 1980 mit der Gegenthese}: Regelwechsel seien selten, und für die laufende Politik genüge die reduzierte Form. Der Streit ist nicht entschieden; beide erhielten den Nobelpreis.}
\end{minipage}
\end{center}

\schritt{Und die Wettbewerbe}
\begin{center}
\begin{minipage}{0.9\textwidth}
\kursiv{\Person{Spyros Makridakis} veranstaltet seit 1982 Vorhersagewettbewerbe mit tausenden Reihen. \textbf{Die ersten drei gewannen einfache Glättungsverfahren}; erst der vierte 2018 wurde von einer Kombination aus statistischen und Lernverfahren gewonnen, der fünfte 2020 von Gradientenbäumen aus Band 22. Die vier Lehren aus Abschnitt 13.3 sind sein Fazit, und sie haben vierzig Jahre gehalten.}
\end{minipage}
\end{center}
\end{geschichte}

\begin{forschung}
\boxtitelp{Vorhersage und Politik heute}
\begin{center}
{\small
\renewcommand{\arraystretch}{1.45}
\begin{tabular}{@{}P{3.2cm}P{5.0cm}@{}}
\toprule
\textbf{Sufficient statistics} & Politikbewertung aus wenigen reduzierten Größen, die auch unter Regelwechsel gelten.\\
\textbf{Erwartungsdaten} & Erwartungen direkt messen statt sie zu unterstellen.\\
\textbf{Lernverfahren für Makro} & Band 22 auf Zeitreihen, mit allen Fragen dieses Kapitels.\\
\bottomrule
\end{tabular}}
\end{center}

\schritt{Zur ersten Zeile}
\begin{center}
\begin{minipage}{0.88\textwidth}
\kursiv{Manche Politikfragen lassen sich mit wenigen Elastizitäten beantworten, die zwar reduzierte Formen sind, aber unter der betrachteten Politikänderung stabil bleiben --- etwa die Preiselastizität der Nachfrage für die Wohlfahrtswirkung einer kleinen Zuzahlungsänderung. \textbf{\Person{Chetty} hat diesen Ansatz 2009 systematisiert}, und er ist der pragmatische Mittelweg zwischen Lucas und Sims: nicht alle tiefen Parameter, aber auch nicht die ganze reduzierte Form --- nur das, was für die Frage zählt und stabil ist.}
\end{minipage}
\end{center}
\end{forschung}

\begin{merksatz}
\boxtitel{Was Sie jetzt können}
\begin{center}
\renewcommand{\arraystretch}{1.5}
\begin{tabular}{@{}P{3.0cm}P{5.6cm}@{}}
\toprule
\textbf{Zwei Zwecke} & vorhersagen gegen eingreifen\\
\textbf{Satz 13.1} & Koeffizienten hängen von der Regel ab\\
\textbf{Tief und reduziert} & $\beta$ gegen $\beta+\gamma(\phi-1)$\\
\textbf{Der Wettbewerb} & klein schlägt groß, Mittel schlägt alle\\
\textbf{Der Bruchfall} & bestes Vorhersagemodell, falsche Politik\\
\bottomrule
\end{tabular}
\end{center}

\begin{center}
\begin{minipage}{0.88\textwidth}
\kursiv{Dieses Kapitel hat gezeigt, was Modelle nicht können. Das nächste kehrt zum Leitfall zurück und fragt, was man mit allen Werkzeugen des Bandes über die Zuzahlung tatsächlich sagen kann --- und wie man es sagt.}
\end{minipage}
\end{center}
\end{merksatz}
